\section{Conclusion}
LEAF separates contextual entity processing within a decision from temporal recurrence across decisions. A shared loop develops receiver-relative entity context, learned fusion combines intermediate reads, and a downstream GRU maintains memory. Query-history input is zero and explicit count injection is disabled. Frozen LEAF policies achieve higher group mean reward $R=-D$ than five learning baselines and Rule across the four evaluated HAD population groups. The saved memory-query comparison favors LEAF at the group-mean level while both designs retain downstream recurrence.

On the three-configuration diagnostic subset, reading the initial workspace without contextual updates lowers all nine seed--configuration reward means despite retaining the relation projection and GRU. The first two rounds receive 98.59--98.91\% of average fusion weight, but most two-round action changes occur among low-margin decisions. Fixing empirical fusion weights has mixed reward effects at $K=2$ and lowers all three seed means at $(30,30,12)$; state-dependent weighting has a consistent advantage in that tested target-rich configuration. GRU reset lowers eight of nine means. These results identify conditional reliance of the saved policies on contextual features, read selection, and temporal state, rather than separate gains from retraining each component. Depth effects remain mixed, and the local action-stability condition does not guarantee closed-loop reward. The evidence concerns a fixed simulated deployment family; transfer across adversary strategies and physical operating conditions remains to be established.
